\section{Reciprocal traces, phase recovery, and height saturation}
\label{sec:reciprocal-trace-height}

This section separates two reciprocal functions which have different
arithmetic meanings.  Put
\begin{equation}\label{eq:Psi-and-Fgamma}
 \Psi(z)=\Ein(z)+\Ein(z^{-1})=H(J(z)),
 \qquad
 \mathcal F_\gamma(z)=E_1(z)+E_1(z^{-1})=\Psi(z)-2\gamma,
\end{equation}
where branches are chosen coherently so that
$\Log(z^{-1})=-\Log z$.  The first function has rational Laurent
coefficients and is genuinely independent of $\gamma$.  The second is the
Euler-level translate of the first; the opposite monodromies of its two
$E_1$ terms cancel, so the displayed sum denotes its single-valued analytic
continuation on $\C^\times$.

This distinction is also essential for effectivity.  The entire function
$\Ein$ is a Siegel $E$-function, whereas $E_1$ has a logarithmic
singularity at the origin and is not an $E$-function.  In particular,
an algebraic approximation to a zero of $\mathcal F_\gamma$ does not put all
terms of
\[
 \Log z+E_1(z)+\Ein(z^{-1})
\]
inside the effective theory of $E$-values.  Indeed the connection formula
gives the identity
\begin{equation}\label{eq:tautological-recovery}
 \Log z+E_1(z)+\Ein(z^{-1})
 =\Psi(z)-\gamma=\gamma+\mathcal F_\gamma(z).
\end{equation}
Thus evaluation at a zero of $\mathcal F_\gamma$ recovers $\gamma$, but
only tautologically.  The nontrivial recovery below uses instead the zero
divisor and the asymptotic phase of reciprocal functions.

\subsection{An exact Hadamard trace at the Euler level}

Let $\{u_k\}$ be the zeros of $H(u)-2\gamma$, repeated according to
multiplicity.  The index set includes the unique real zero from
Theorem~\ref{thm:H-euler-level} and both members of every nonreal conjugate
pair.

\begin{theorem}[Euler-level Hadamard trace]
\label{thm:euler-hadamard-trace}
The series $\sum_k |u_k|^{-2}$ converges and
\begin{equation}\label{eq:euler-hadamard-trace}
 \boxed{
 \sum_k\frac{1}{u_k^2}
 =\frac{\sin^2 1-\sin 1\,\Ci(1)}{4\Ci(1)^2}.}
\end{equation}
Consequently, if $T$ denotes the left side, then
\begin{equation}\label{eq:trace-inversion}
 \Ci(1)=
 \frac{\sin 1\bigl(\sqrt{1+16T}-1\bigr)}{8T},
 \qquad
 \gamma=\Cin(1)+
 \frac{\sin 1\bigl(\sqrt{1+16T}-1\bigr)}{8T}.
\end{equation}
\end{theorem}

\begin{proof}
Set $f(u)=H(u)-2\gamma$.  By
Corollary~\ref{cor:H-order-zero-count}, $f$ is entire of order one.  Hence
its zero divisor has exponent of convergence at most one and admits a
genus-one Hadamard factorization.  Since
\[
 f(0)=2\Cin(1)-2\gamma=-2\Ci(1)\ne0,
\]
logarithmic differentiation twice at the origin gives
\begin{equation}\label{eq:hadamard-log-derivative}
 -\left(\frac{f'}f\right)'(0)=\sum_k\frac1{u_k^2}.
\end{equation}
The product formula for $H'$ in
Proposition~\ref{prop:entire-joukowski-descent} gives the exact origin jet
\[
 H'(0)=\sin 1,
 \qquad
 H''(0)=-\frac12\sin 1.
\]
Therefore
\[
 -\left(\frac{f'}f\right)'(0)
 =\frac{f'(0)^2-f''(0)f(0)}{f(0)^2}
 =\frac{\sin^2 1-\sin 1\,\Ci(1)}{4\Ci(1)^2},
\]
which proves \eqref{eq:euler-hadamard-trace}.  The standard bounds
$0<\Ci(1)<\sin 1$ (which also follow directly from the corresponding
alternating series and elementary bounds for $\gamma$) show that the trace
$T$ is positive.  Solving
$4T X^2+\sin(1)X-\sin^2(1)=0$ for the positive root
$X=\Ci(1)$ gives \eqref{eq:trace-inversion}.
\end{proof}

\begin{remark}[What the trace does not prove]
Formula \eqref{eq:trace-inversion} is an exact spectral reparametrization,
not a transcendence proof: the divisor on the left is itself the divisor of
the $\gamma$-dependent level $H-2\gamma$.  No arithmetic assertion in this
paper relies on a numerical zero count.  In particular, an exact floor
formula suggested by finite computations is not used; the rigorous bounds
remain those of Corollary~\ref{cor:H-order-zero-count}.
\end{remark}

\subsection{A nontrivial phase recovery from the
\texorpdfstring{$\gamma$}{gamma}-free zero set}

The situation is different for zeros of $\Psi$, because the defining
Laurent series then contains no $\gamma$.  Let $z_{k,0}$ be the upper
half-plane zero with $|z_{k,0}|>1$ constructed in the proof of
Theorem~\ref{thm:H-level-asymptotics} at level $c=0$, and write
$\rho_k=z_{k,0}$.  Put
\[
 t_k=2\pi k,
 \qquad q_k=\log t_k.
\]

\begin{theorem}[Second-order zero phase]
\label{thm:gamma-free-zero-phase}
As $k\to\infty$,
\begin{align}
 \rho_k={}&-q_k-\frac12\log\left((q_k+\gamma)^2+\frac{\pi^2}{4}\right)
 \label{eq:rho-second-order}\\
 &\quad+i\left[t_k+\arctan\left(\frac{2(q_k+\gamma)}{\pi}\right)\right]
 +O\left(\frac{q_k}{t_k}\right).\nonumber
\end{align}
In particular the zero set of the $\gamma$-free function $\Psi$ determines
$\gamma$ through
\begin{equation}\label{eq:gamma-phase-recovery}
 \boxed{
 \gamma=\lim_{k\to\infty}
 \left\{
  \frac\pi2\tan\bigl(\Im\rho_k-2\pi k\bigr)-\log(2\pi k)
 \right\}.}
\end{equation}
More precisely, the expression in braces is
\begin{equation}\label{eq:gamma-phase-rate}
 \gamma+O\!\left(\frac{(\log k)^3}{k}\right).
\end{equation}
Every $\rho_k$ is transcendental.
\end{theorem}

\begin{proof}
At level $c=0$, the fixed-point equation in the proof of
Theorem~\ref{thm:H-level-asymptotics} is
\begin{equation}\label{eq:refined-fixed-point}
 z=(2k+1)\pi i-\Log z-\Log(\Log z+\gamma).
\end{equation}
The exact zero $\rho_k$ differs from its model fixed point by $O(t_k^{-1})$.
Uniformly in the unit disc used there,
\[
 \Log z=q_k+\frac{\pi i}{2}+O(q_k/t_k),
\]
and hence
\begin{align*}
 \Log(\Log z+\gamma)
  ={}&\frac12\log\left((q_k+\gamma)^2+\frac{\pi^2}{4}\right)\\
   &+i\arctan\left(\frac{\pi}{2(q_k+\gamma)}\right)
     +O(q_k/t_k).
\end{align*}
Substitution in \eqref{eq:refined-fixed-point}, followed by
\[
 \frac\pi2-\arctan\left(\frac{\pi}{2x}\right)
 =\arctan\left(\frac{2x}{\pi}\right)
 \qquad(x>0),
\]
proves \eqref{eq:rho-second-order}.

Write
$\Im\rho_k-t_k=\arctan(2(q_k+\gamma)/\pi)+\varepsilon_k$, where
$\varepsilon_k=O(q_k/t_k)$.  Since $q_k/t_k=o(q_k^{-1})$, the whole
mean-value segment remains a distance asymptotic to $q_k^{-1}$ from the
nearest pole of $\tan$.  Consequently $\sec^2=O(q_k^2)$ uniformly there,
\[
 \frac\pi2\tan(\Im\rho_k-t_k)
 =q_k+\gamma+O(q_k^3/t_k).
\]
The error tends to zero and proves both
\eqref{eq:gamma-phase-recovery} and \eqref{eq:gamma-phase-rate}.
Finally $\Psi(\rho_k)=0$, so $\rho_k$ lies over the algebraic level
$c=0$; its transcendence follows from
Corollary~\ref{cor:H-algebraic-levels}.
\end{proof}

Theorem~\ref{thm:gamma-free-zero-phase} is analytically nonvacuous: unlike
\eqref{eq:tautological-recovery}, the object whose zeros are used has
rational Laurent coefficients.  It still gives no irrationality result,
because no effective arithmetic control of these transcendental zeros is
established here.

\subsection{Laurent truncations saturate the height budget}

Define
\begin{equation}\label{eq:Psi-truncation}
 \Psi_N(z)=\sum_{m=1}^N
 \frac{(-1)^{m+1}}{m\,m!}(z^m+z^{-m}),
 \qquad
 L_N=\operatorname{lcm}(1,\ldots,N),
\end{equation}
and
\begin{equation}\label{eq:QN-definition}
 D_N^*=\operatorname{lcm}_{1\le m\le N}(m\,m!),
 \qquad
 Q_N(z)=D_N^*z^N\Psi_N(z).
\end{equation}

\begin{lemma}[The denominator scale is intrinsic]
\label{lem:truncation-denominator-scale}
One has
\[
 Q_N\in\Z[z],\qquad Q_N\ \text{primitive},\qquad \deg Q_N=2N,
 \qquad H(Q_N)=D_N^*,
\]
and therefore
\begin{equation}\label{eq:QN-height}
 \log H(Q_N)=N\log N+o(N).
\end{equation}
Moreover,
\begin{equation}\label{eq:minimal-denominator-sandwich}
 N!\prod_{p\le N}p\ \mid\ D_N^*\ \mid\ N!L_N,
\end{equation}
so $\log D_N^*=N\log N+o(N)$ as well.
\end{lemma}

\begin{proof}
The coefficients in \eqref{eq:QN-definition} are
$D_N^*/(m\,m!)$, hence are integral and have maximum absolute value
$D_N^*$, attained at $m=1$.  Their gcd is one: otherwise division by a
common prime would give a smaller common denominator than $D_N^*$.
This proves the first assertions.

For the right divisibility in \eqref{eq:minimal-denominator-sandwich}, use
$m!\mid N!$ and $m\mid L_N$.  For the left divisibility, fix a prime
$p\le N$ and take $m=p\lfloor N/p\rfloor$.  There is no multiple of $p$
strictly between $m$ and $N$, so
$v_p(m!)=v_p(N!)$ and
$v_p(m\,m!)\ge v_p(N!)+1$.  Taking the least common multiple over $m$
adds at least one copy of every prime $p\le N$ to $N!$.  The prime number
theorem and Stirling's formula give \eqref{eq:QN-height}.
\end{proof}

\begin{theorem}[Factorial approximation and defining-height saturation]
\label{thm:truncation-height-saturation}
Let $\rho$ be a fixed simple zero of $\Psi$ with $|\rho|>1$.  For all
sufficiently large $N$, $\Psi_N$ has a zero $\alpha_N$ converging to
$\rho$, and
\begin{equation}\label{eq:truncation-root-asymptotic}
 \boxed{
 \alpha_N-\rho=
 \frac{(-1)^{N+2}\rho^{N+1}}
 {\Psi'(\rho)(N+1)(N+1)!}
 \left(1+O_\rho(N^{-1})\right).}
\end{equation}
Consequently
\begin{equation}\label{eq:defining-height-ratio}
 \frac{-\log|\alpha_N-\rho|}{\log H(Q_N)}\longrightarrow1.
\end{equation}
Here $Q_N$ is the primitive integer polynomial obtained by minimal
denominator clearing; no claim about the minimal-polynomial height of
$\alpha_N$ is made.
\end{theorem}

\begin{proof}
On a fixed neighbourhood of $\rho$, the tail
$R_N=\Psi-\Psi_N$ satisfies
\begin{equation}\label{eq:Psi-tail-leading}
 R_N(z)=
 \frac{(-1)^{N+2}z^{N+1}}{(N+1)(N+1)!}
 \left(1+O_\rho(N^{-1})\right).
\end{equation}
Indeed the ratio of consecutive exterior terms is $O_\rho(N^{-1})$, while
the terms in $z^{-m}$ are exponentially smaller.  Rouch\'e's theorem on a
small circle about the simple zero $\rho$ gives a unique nearby zero
$\alpha_N$ of $\Psi_N$.  The standard perturbation formula yields
\[
 \alpha_N-\rho
 =\frac{R_N(\rho)}{\Psi'(\rho)}
  \left(1+O_\rho(N^{-1})\right),
\]
which is \eqref{eq:truncation-root-asymptotic}.  Stirling's formula and
Lemma~\ref{lem:truncation-denominator-scale} give
\[
 -\log|\alpha_N-\rho|=N\log N+O_\rho(N),
 \qquad
 \log H(Q_N)=N\log N+o(N),
\]
and hence \eqref{eq:defining-height-ratio}.
\end{proof}

\begin{corollary}[Full primitive integerization destroys the small value]
\label{cor:integerized-truncation-growth}
Under the hypotheses of
Theorem~\ref{thm:truncation-height-saturation},
\begin{equation}\label{eq:cleared-value-asymptotic}
 \boxed{
 Q_N(\rho)=
 (-1)^{N+3}\frac{D_N^*}{(N+1)(N+1)!}\rho^{2N+1}
 \left(1+O_\rho(N^{-1})\right).}
\end{equation}
In particular
\begin{equation}\label{eq:cleared-value-growth}
 \log|Q_N(\rho)|
 =N\bigl(1+2\log|\rho|\bigr)+o(N)>0.
\end{equation}
Also
\begin{equation}\label{eq:moving-point-value-scale}
 |\Psi(\alpha_N)|=H(Q_N)^{-1+o(1)}.
\end{equation}
\end{corollary}

\begin{proof}
Since $\Psi(\rho)=0$, one has
$Q_N(\rho)=-D_N^*\rho^N R_N(\rho)$.  Substitution of
\eqref{eq:Psi-tail-leading} gives \eqref{eq:cleared-value-asymptotic}.
Lemma~\ref{lem:truncation-denominator-scale} and Stirling's formula then
give \eqref{eq:cleared-value-growth}.  Finally
$\Psi_N(\alpha_N)=0$, so the same tail estimate at $\alpha_N$ and
\eqref{eq:QN-height} give \eqref{eq:moving-point-value-scale}.
\end{proof}

The last corollary closes the naive use of this full primitive clearing
polynomial.  Taken alone, it does not exclude systematic factorization or a
minimal-height compression for the selected root; the dyadic theorem in
Section~\ref{sec:dyadic-capacity} closes precisely that escape on
\(N=2^r\).  The factorial analytic
error is real, but after denominator clearing the full primitive polynomial
value grows rather than tends to zero.  It also explains
why the pointwise measure of Adamczewski--Faverjon cannot simply be compared
with $H(Q_N)$: in their theorem the height belongs to the relation
polynomial in fixed $E$-values, whereas here the algebraic evaluation point
and its degree move with $N$.

\subsection{The exterior-rank barrier}

Let $C_m$ be the monic Chebyshev polynomials from
\eqref{eq:joukowski-chebyshev} and set
\begin{equation}\label{eq:HN-truncation}
 H_N(w)=\sum_{m=1}^N\frac{(-1)^{m+1}}{m\,m!}C_m(w).
\end{equation}

\begin{theorem}[The level first appears in the top norm]
\label{thm:top-norm-barrier}
Fix a level $\ell\in\C$ and let
$w_{1,N},\ldots,w_{N,N}$ be the roots, with multiplicity, of
$H_N(w)-\ell$.  Then the elementary symmetric functions
\[
 e_1(w_{1,N},\ldots,w_{N,N}),\ldots,
 e_{N-1}(w_{1,N},\ldots,w_{N,N})
\]
are independent of $\ell$, while
\begin{equation}\label{eq:top-norm-exact}
 \boxed{
 \prod_{j=1}^N w_{j,N}
 =N\,N!\bigl(\ell-H_N(0)\bigr).}
\end{equation}
Thus the first elementary symmetric invariant which sees the level is the full
rank-$N$ norm, and it carries the factorial factor $N!$.
\end{theorem}

\begin{proof}
Since $C_N$ is monic, the leading coefficient of $H_N-\ell$ is
$(-1)^{N+1}/(N\,N!)$.  The level $\ell$ changes only the constant
coefficient.  Vieta's formulas therefore make $e_1,\ldots,e_{N-1}$
level-independent, and give
\[
 e_N=(-1)^N
 \frac{H_N(0)-\ell}{(-1)^{N+1}/(N\,N!)}
 =N\,N!\bigl(\ell-H_N(0)\bigr).
\]
\end{proof}

At the Euler level $\ell=2\gamma$, since
$H_N(0)\to H(0)=2\Cin(1)$, one has
\[
 \prod_{j=1}^Nw_{j,N}=2\Ci(1)\,N\,N!\,(1+o(1)),
 \qquad
 \log\left|\prod_{j=1}^Nw_{j,N}\right|
 =N\log N-N+O(\log N).
\]
Thus, by Newton's identities, the power sums of order below $N$ and all
proper exterior traces of the finite zero set forget $\gamma$; the first
elementary symmetric invariant which
remembers it pays the full determinant cost.  A possible route not excluded
by the theorem
would have to remove the universal $N\log N$ contribution canonically, for
example through a zeta-regularized determinant or an arithmetic spectral
projector.  No such arithmetic regularization is constructed here.
